教育界流行着一句名言：要把金针度与人。这儿的关键是“金针”二字，而不能是“铁针”，更不能是“锈针”。语文教学的“金针”是什么呢？我认为，应指带有规律性的语文基本知识和运用这些规律的方法。教师如果在教学生如何在把握住“金针”、训练学生具有自学能力方面多下功夫，就能使教学内容少而精，同时在教学方法上必然要以启发式代替填鸭式；学生如果在学习如何把握好“金针”、提高自学能力方面多下功夫，就能充分发挥学习语文的主观能动性，逐渐做到触类旁通，从而收到事半功倍的效果。

为了探索这些规律和方法，我曾写过十篇短文，承蒙《语文报》编辑同志将它们以“语言小钥匙”、“写作谈片”等专栏的名目予以发表。我的同事和学生看了这些短文之后，认为不论是对教语文、还是对学语文都有所帮助，希望并鼓励我能就语文能力训练的各个方面、各个环节，作比较系统的探讨。就在同志们的关心和鼓励下，我在去年初完成了《中学语文全程训练》的书稿，并交付出版。

今年四月，接到陕西人民教育出版社编辑同志来信，谈到《中学语文全程训练》一书在出版一个月后即已脱销，而各地读者还纷纷去函订购此书。鉴于这种情况，出版社编辑同志建议我修订此书，交予付印。我再三考虑，认为读者最感兴趣的，恐怕是如何去掌握打开语文大门的“钥匙”。于是，我花了几个月的功夫，将全书近乎重写一道，不仅彻底删去原书的练习和答案，而且对原书的四十个论题逐一地进行了较大幅度的修改；另外，为了适应读者的需要，弥补原书的不足，我还特地新写了十个论题，增加了两个“附录”。这样，就把本书的重点由全程训练转到了能力培养方面来，突出了培养举一反三的自学能力这个为笔者所孜孜以求的目标。正因如此，特将书名题为《语文开窍》。

本书在探讨培养语文自学能力的道路上，仅只是一块铺路砖，它还很粗糙。如果读者感到它能对如何学得有方这个问题的解决起到一定的作用，我就很满意了。

在编写过程中，我吸取了很多有关这方面的研究成果，在此一并谨致谢意。

陕西人民教育出版社的编辑同志对于此书的修改和定稿提出许多宝贵的建议，花费大量的劳动，也请接受我由衷的敬意和谢意。
\begin{flushright}
	笔\hspace{1em}者\hspace{3em}
	
	一九八六年九月五日\hspace{4em}
	
	于连云港教育学院
\end{flushright}
